%--------------------------------------------------------------
\section{先行研究との比較}\label{sec:compare}
%--------------------------------------------------------------
\subsection{Hodge 理論（形式係数）}
Gualtieri は一般化 Kähler 多様体上の一般化 Hodge star、Born--Infeld 内積、Laplacian の等式、Hodge 分解、$dd^{\Jc}$-補題を示した（\tC\ Gualtieri, \emph{Generalized geometry and the Hodge decomposition}, arXiv:math/0409093；Gualtieri, \emph{Generalized K\"ahler geometry}, Comm.\ Math.\ Phys.\ 331 (2014) 297--331, arXiv:1007.3485；Cavalcanti, \emph{New aspects of the $dd^c$-lemma}, arXiv:math/0501406、および J.\ Geom.\ Phys.\ 57 (2006) 121--132。定理番号は精読していないので省く\tQ）。これらは自明束または $H$-捩れ係数の理論であり、直線束（曲率つき）係数の Bochner--Kodaira--Nakano 恒等式は、調べた範囲では書かれていない。本ノートの結果は、この係数付き版の穴を、Hodge 理論に頼らず、対数複体と構成可能層の比較で $\beta$-変換の場合に埋めたものと位置づけられる。

\subsection{Pal--Poddar}
\tC\ Pal--Poddar, \emph{Strong generalized holomorphic principal bundles}, arXiv:2404.18113 は、正則型（型 $k$ が一定）のコンパクト一般化複素多様体で、シンプレクティック葉構造の葉空間が Kähler 軌道体である場合に、強一般化正則（SGH）束の Hodge 理論、横断 Dolbeault 同型（Thm.\ 8.9, Cor.\ 8.10。番号は\tQ）、Hodge 分解と Serre 双対性（Thm.\ 11.2, 11.4, 11.6, 11.7）、Kähler の恒等式（Prop.\ 11.8, Lemma 11.9。番号は\tQ）、正の SGH 直線束（Def.\ 11.10, 11.11）に対する Kodaira 型・Serre 型の消滅定理（Thm.\ 11.12）を与える。正則型と軌道体の仮定は、葉構造が滑らかな葉束をなすことと、横断方向の Hodge 理論（楕円性）を保証するために使われている。本ノートとの違い：
\begin{itemize}
\item 型が跳ぶ場合（$\sigma$ の退化因子上で型が跳ぶ）は Pal--Poddar の枠外であり、本ノートの $\mathbb{CP}^2$ などの例はこの場合である。
\item SGH 束は一般化正則束の真部分クラス（葉に沿って平坦で横断方向に正則）。本ノートの Poisson 加群は $\beta$-変換構造における一般の GH 束である。
\item 定理 11.12 の主張・証明の構造は \S11 を再読して確認した。正確な主張：$H^q(M,(\mathcal G^*M)^p\otimes E)=0$ $(p+q>k)$、および $H^q(M,E'\otimes E^m)=0$ $(m\ge m_0,\ q>0)$。仮定（コンパクト、葉空間が Kähler 軌道体、$E$ が正）は、横断方向の Hodge 理論と Bochner--Kodaira--Nakano 型の曲率項の正値性に使われる。
\end{itemize}
原典の記述について、本調査で気づいた疑義（原典の誤植・省略の可能性が高いが未検証\tQ）：(i) 主張は $(i/2\pi)$ 付き、証明では $1/2\pi$ なし（無害）、(ii) $gc_1$ の符号の食い違いの原因は、特性多項式 (9.9) と \S10.1 の $\det(I\pm tA/2\pi i)$ の定義の食い違い、(iii)「正のクラスから正曲率の計量が出る」ことは Def.\ 11.11 の定義から証明されず、Thm.\ 11.12 の証明が仮定している、(iv) Def.\ 11.11 にコンパクト性・Kähler 性の明記がない、(v) Lemma 11.9(1) の証明は点ごとの局所枠での計算、(vi) Thm.\ 13.4(2)(3) の $T^{1,0}$ 対 $T^{1,0*}$、Example 13.10 の $q=n$ の二項係数。

\subsection{Tengzhou Hu の仕事}
\tC\ T.\ Hu, \emph{Kodaira vanishing theorems for K\"ahler Lie algebroids}, arXiv:2403.10876；\\ T.\ Hu, \emph{Lie groupoid Riemann--Roch--Hirzebruch theorem and applications}, arXiv:2403.13424。
\begin{itemize}
\item \emph{設定}（2403.13424）：Lie 亜群 $G\rightrightarrows M$（Def.\ 1、固有性・$s$-連結性の仮定は述べられていない）、その Lie 亜群 $\mathfrak g$。「複素 Lie 亜群」（Def.\ 2）は、\emph{実}の Lie 亜群 $\mathfrak g$ に、ファイバー方向の可積分な複素構造 $J$ を与えたもの（これは $\bT_\C$ の部分束としての複素 Lie 亜群 $L$ とは別物）。Kähler Lie 亜群は Def.\ 11、正の束は Def.\ 12。
\item \emph{主張}（2403.13424）：Thm.\ 4（Riemann--Roch）は、$M$ コンパクト、$\dbar_{\mathfrak g}$ が楕円型（主張されているが証明は見当たらない）、$\Omega$ が $\mathfrak g$-不変な密度（単模性）のとき
\[\mathrm{Ind}_\Omega(\dbar_{\mathfrak g}+\dbar_{\mathfrak g}^*)=(2\pi i)^{-k}\int_M(-1)^{\mathrm{rk}\,\mathfrak g}\bigl\langle\ch^{\mathfrak g}(\mathfrak g^{*p,0})\ch^{\mathfrak g}(E)\Td^{\mathfrak g}_\C(\mathfrak g),\Omega_{\mathfrak g}\bigr\rangle\]
を主張する。Thm.\ 7（消滅）：$E$ が正で $p+q>n$ なら $\mathfrak g$-Dolbeault Laplacian の\emph{核}が消える（コホモロジーの消滅は導けない：$\mathfrak g$ が一般の Lie 亜群のとき Laplacian は楕円型とは限らず、Hodge 理論が使えないため）。2403.10876 では Kähler Lie 亜群は Def.\ 2、Kähler の恒等式は Thm.\ 1 と Prop.\ 4、BKN 恒等式は Lemma 4、主な消滅定理は序論。
\item \emph{一般化 Dolbeault 複体は適用範囲外}\tP\ （導出、原典との照合は以下の範囲）：$L$ は複素 Lie 亜群で $J$ をもたず、$\wedge L^*$ は一般には $(p,q)$ 二重次数をもたない。正則 Poisson 構造 $\sigma$ に対し、実 Lie 亜群 $A_{\mathbb R}$ は Hu の設定に入るが、その Dolbeault 複体は $\bar\sigma$ から作られる $\wedge\mathfrak g^{*0,q}$ で、本ノートの複体 $(T^{0,1}\bowtie A^{1,0})$ とは別物。
\item \emph{指数は $m$ を見るか}\tK：Hu の指数は群の指数で、おおよそ $\int\ch(E)\Td(T^*)\ch(\wedge^pT)$ のように $m$ に多項式的に依存する見込みだが、本ノートの全 Euler 標数は $m$ によらない（命題 \ref{prop:euler}）。また周期的な跳びを決める留数 $-m/3$ は Hu のデータに含まれない。したがって Hu の定理は「一般化 Dolbeault コホモロジーに対する Riemann--Roch」ではない。
\item \emph{原典の不整合}：2403.13424（要約を通して得たもの）：Thm.\ 7 で $n$ が未定義、符号定数が Thm.\ 5, Prop.\ 1, Thm.\ 6 で異なる、「単模」の定義が二通り。2403.10876（本文の続きと付録の一部を読んで確認\tC）：Lemma 5 は $[L,\Lambda]=(p+q-n)\mathrm{id}$ だが、Thm.\ 2 の証明は $\Delta_{\bar\partial_{\mathfrak g}}=\Delta^{1,0}-[L,\Lambda]=\Delta^{1,0}+(p+q-n)\mathrm{id}$ と書き、後の等号の符号が合わない（$-(p+q-n)=n-p-q$）。さらに代入 $\nabla^2=\sqrt{-1}L$ は $\sqrt{-1}R=-L$ を意味し、「正」（$\sqrt{-1}R>0$、Def.\ 13）の定義と逆である。主張（正の $E$ で $p+q>n$ のとき核が消える）自体は古典の Kodaira--Nakano と整合する。つまり印刷された証明に符号のずれがある（主張の否定ではない）。Cor.\ 1 に矢印と記号 $s/t$ の誤植。付録（Lemma 6--10、$[\bar\partial_{\mathfrak g},\iota_\omega]$ の計算）は添字が読み取りにくく検証できていない。参考文献欄と掲載誌は未確認。\tQ
\item \emph{本ノートの前版（予備調査）の訂正}：前版では著者を「S.\ Hu」、題名を別の題名で書いていたが、いずれも誤りで、上記が正しい。
\end{itemize}

\subsection{積分可能性・群}
\tC\ Gualtieri--Li（arXiv:1206.3674、IMRN 2014）は対数シンプレクティック多様体の Hausdorff シンプレクティック亜群を分類し、滑らかな $D$ に沿って接する対数ベクトル場の Lie 亜群 $T_X(-\log D)$ が積分可能であることを与える。Bailey--Gualtieri（arXiv:1611.03850）は、一般化複素構造の弱正則シンプレクティック亜群への積分と実 Poisson 構造の積分可能性の同値を示す（Thm.\ 7.1：下の実 Poisson 構造が積分可能であることが必要十分。Thm.\ 7.1 自体は Hausdorff・$s$-連結性を述べない。構成 Thm.\ 7.6 は $s$-連結かつ $s$-単連結な積分を用いる）。Laurent-Gengoux--Stiénon--Xu（arXiv:0803.2031）：Thm.\ 3.17, 3.22。Crainic--Fernandes（arXiv:math/0105033）：Thm.\ 4.1, Cor.\ 5.2。これらは、指数理論の入力としての群を与えるが、複素 Lie 亜群 $L$ 自体の積分は扱わない。

%--------------------------------------------------------------
\section{未解決問題と次のステップ}\label{sec:open}
%--------------------------------------------------------------
\begin{enumerate}
\item \textbf{三種の神器のうち未着手のもの}：Kodaira 埋め込みの一般化。$\HH^0$ が Casimir 切断であること（観察 \ref{obs:casimir}）から、$M$ 自身への埋め込みではなく、Casimir 切断が分離する商（葉空間など）への写像として定式化する必要がある。Casimir 切断の線形系のベースポイント・分離性を、$\mathbb{CP}^2$ の例（$\HH^0=\C F^{m/3}$、すなわち $m\in3\Z_{\ge0}$ で 1 次元）から調べるのが次の具体的な一歩。
\item \textbf{Serre 型消滅の一般化}（予想 \ref{conj:A3}）：「正」な一般化正則直線束の定義、$\beta$-変換でない $\Jc$ での $E_1$ 項に相当する濾過の構成。Pal--Poddar Thm.\ 11.12 との比較は \S11 の再読で行った（群は別物、横断次数の読みは整合）。
\item \textbf{Riemann--Roch の定式化}：全 Euler 標数ではなく、細分（$\Jc$ の型による濾過の $E_1$ 項、周期擬多項式、指標根 $H^*(\partial U)$）を決める。予想 \ref{conj:I3} の証明には、複素 Lie 亜群 $L$ の指数理論（単模性・Fredholm 性 $\Leftrightarrow r\notin\Z$ の導出）が必要。境界つき一般化複素多様体への APS 型の拡張は、さらに先の課題。
\item \textbf{対数シンプレクティック以外}：非ノーダル特異点の局所理論の証明（準斉次 germ のみ計算）、留数が等しくない場合の大域公式、非簡約 $D$（無限次元の茎と線形増大）、次元 4 以上の非トーリックな例（観察 \ref{obs:arith} を参照）。
\item \textbf{一般化 Kähler 側}：直線束係数の Bochner--Kodaira--Nakano 恒等式を一般化 Kähler の枠組みで書き下す。$H$-捩れと曲率項の相互作用が障害になりうる\tK。
\item \textbf{本ノートの未確認事項の解消}：\S\ref{sec:ledger} の \tQ\ の項目。特に、Mokri の定理、LGSX Lemma 4.22 の再導出、Gualtieri 2011 \S4 の精読、残る書誌情報。
\end{enumerate}

\paragraph{次のステップの優先順位（筆者案）。}
(a) 型が跳ばない比較対象として、Pal--Poddar の正則型の例で同じ「Poincaré 多項式」を計算して、本ノートの周期的な跳びと見比べる。(b) $\mathbb{CP}^2$ の例について、$\HH^0=\C F^{m/3}$ を使った Casimir 切断の線形系の幾何を調べ、Kodaira 型定理の「正しい形」を探る。(c) 予想 \ref{conj:I3} の指数理論的な証明の足場として、実 Lie 亜群 $T_X(-\log D)$ 上の重みつき $b$-複体と留数 $r$ の関係を厳密に定式化する。
